\subsection{等价关系的商}

设R、S是关于两个字母x、y的两个等价关系。若关系 \(S \Longrightarrow R\) 成立。若R和S是同一集合E上的等价关系，则称S比R更细意味着S的图包含于R的图中，或者说S的每个等价类都包含于R的某个等价类中；等价地，R的每个等价类关于S都是饱和的。

示例

\begin{enumerate}
\def\labelenumi{(\arabic{enumi})}
\tightlist
\item
  关系“ \(x \in E\) 并且 \(y \in E\) 并且 \(x = y\) “”比E上的每一个等价关系都更细。关系“ \(x \in E\) 并且 \(y \in E\) “” 比 E 上的每一个等价关系都更粗。
\end{enumerate}

* (2) 等价关系“ \(x \in \mathbf{Z}\) 并且 \(y \in \mathbf{Z}\) 并且 \(x - y\) “能被4整除”比等价关系“ \(x \in \mathbf{Z}\) 并且 \(y \in \mathbf{Z}\) 并且 \(x - y\) 能被2整除''。*

设R和S是同一集合E上的两个等价关系，且S比R更细。设f和g分别是E到E/R和E到E/S的典范映射；则函数f与S相容。设h是f在关于S取商后诱导出的函数；则h是E/S到E/R的一个映射。与h在E/S上相关联的等价关系称为R关于S的商，记为R/S。该关系 \(x \equiv y \pmod{R}\) 等价于 \(g(x) \equiv g(y) \pmod{R/S}\) ，并且关于 R/S 的等价类正是 E 中关于 R 的等价类在 g 下的像。设 \(h = j \circ h_{2} \circ h_{1}\) 为映射 h 的典范分解（第 5 节）。则 \(h_{1}\) 是 E/S 到 (E/S)/(R/S) 上的典范映射，j 是 E/R 的恒等映射，而 \(h_{2}\) 是 (E/S)/(R/S) 到 E/R 上的一一映射。映射 \(h_{2}\) 及其逆映射称为典范映射。

反之，考虑集合 E/S 上的任意等价关系 T，并设 R 为 E 上在 g 下关系 T 的逆像（第 6 节）。由于关系 \(x \equiv y \pmod{R}\) 等价于 \(g(x) \equiv g(y) \pmod{T}\) ，可知 T 等价于 R/S。

